\begin{table}
    \footnotesize
    \centering
    \begin{threeparttable}
        \caption{LLM tonal characteristics by gender}
        \label{table1llm}
        \begin{tabular}{p{5cm}S@{}S@{}S@{}}
            \toprule
            &{Men}&{Women}&{Difference}\\
            \midrule
            \mrow{5cm}{LLM Readability}   &        6.73&        7.39&        0.67***\\
                                          &      (0.02)&      (0.05)&      (0.05)   \\
            \mrow{5cm}{Modal Verb Strength}&        6.17&        6.10&       -0.07   \\
                                          &      (0.03)&      (0.10)&      (0.11)   \\
            \mrow{5cm}{Hedging Frequency \& Type}&        8.35&        8.34&       -0.02   \\
                                          &      (0.02)&      (0.07)&      (0.07)   \\
            \mrow{5cm}{Qualifier Density} &        8.24&        8.30&        0.06   \\
                                          &      (0.02)&      (0.06)&      (0.07)   \\
            \mrow{5cm}{Acknowledgement of Limitations}&        8.05&        8.46&        0.41***\\
                                          &      (0.03)&      (0.09)&      (0.09)   \\
            \mrow{5cm}{Caution-Signaling Connectors}&        8.61&        8.54&       -0.08   \\
                                          &      (0.02)&      (0.08)&      (0.08)   \\
            \mrow{5cm}{Assertiveness \& Voice}&        8.20&        8.27&        0.07   \\
                                          &      (0.02)&      (0.05)&      (0.06)   \\
            \mrow{5cm}{Active/Passive Voice Ratio}&        7.22&        8.24&        1.02***\\
                                          &      (0.03)&      (0.09)&      (0.10)   \\
            \mrow{5cm}{Sentence Length \& Directness}&        5.97&        6.63&        0.66***\\
                                          &      (0.02)&      (0.07)&      (0.07)   \\
            \mrow{5cm}{Imperative-Form Occurrence}&        1.02&        1.00&       -0.02** \\
                                          &      (0.00)&      (0.01)&      (0.01)   \\
            \mrow{5cm}{Pronoun Commitment}&        4.19&        5.40&        1.21***\\
                                          &      (0.04)&      (0.12)&      (0.13)   \\
            \mrow{5cm}{Novelty-Claim Strength}&        2.27&        2.11&       -0.15** \\
                                          &      (0.02)&      (0.07)&      (0.08)   \\
            \mrow{5cm}{Jargon/Technicality Density}&        4.02&        4.69&        0.67***\\
                                          &      (0.02)&      (0.06)&      (0.06)   \\
            \mrow{5cm}{Emotional Valence} &        1.12&        1.11&       -0.01   \\
                                          &      (0.00)&      (0.01)&      (0.01)   \\
            \mrow{5cm}{Evidence \& Citation Usage}&        1.70&        2.24&        0.54***\\
                                          &      (0.02)&      (0.05)&      (0.05)   \\
            \mrow{5cm}{Practical/Impact Orientation}&        2.37&        2.61&        0.23***\\
                                          &      (0.02)&      (0.06)&      (0.06)   \\
            \midrule
            No. observations              &       8,268&         849&       9,117   \\
            \bottomrule
        \end{tabular}
        \begin{tablenotes}
            \tiny
            \item \textit{Notes}. Figures are means of LLM-scored tonal characteristics by sex. Male-authored papers are defined as having a ratio of female authors below 50 percent; female-authored papers are those with a ratio of female authors at or above 50 percent. Last column subtracts male means from female means. Standard errors in parentheses. ***, ** and * difference statistically significant at 1\%, 5\% and 10\%, respectively.
        \end{tablenotes}
    \end{threeparttable}
\end{table}